\documentclass[9pt]{book}
\setlength{\oddsidemargin}{-.5 in}
\setlength{\evensidemargin}{-.5 in}
\addtolength{\topmargin}{-1in}
\setlength{\textwidth}{7.5in}
\setlength{\textheight}{9in}
\usepackage{enumitem}
\usepackage{tikz}

\usepackage{amsmath, amssymb}
\usepackage{algorithm}
\usepackage{url, color, epsfig, amsmath, amssymb}
\usepackage{graphicx}
\usepackage{amsmath, amsthm, amssymb}
\usepackage{algpseudocode}

\usepackage{algorithm}
\usepackage{url, color, epsfig, amsmath, amsthm, amssymb}
\usepackage{amsthm}
\usepackage{graphicx}
\usepackage{mathtools}
\DeclarePairedDelimiter{\ceil}{\lceil}{\rceil}
\DeclarePairedDelimiter{\floor}{\lfloor}{\rfloor}

\newcommand{\remove}[1]{}

\newcommand{\Do}{{\small\bf do}\ }
\newcommand{\Proc}[1]{#1\+}
\newcommand{\Returns}{{\small\bf returns}}
\newcommand{\Procbegin}{{\small\bf begin}}
\newcommand{\Then}{{\small\bf then}\ \=\+}
\newcommand{\Elseif}{\<{\small\bf elseif}\ }
\newcommand{\Endif}{\<{\small\bf end\ if\ }\-\-}
\newcommand{\Endproc}[1]{{\small\bf end} #1\-}
\newcommand{\Endfor}{{\small\bf end\ for}\ \-}
\newcommand{\Endloop}{{\bf end\ loop}\ \-}
\newenvironment{program}{
	\begin{minipage}{\textwidth}
		\begin{tabbing}
			\ \ \ \ \=\kill
		}{
	\end{tabbing}
\end{minipage}
}

\def\blankline{\hbox{}}

\newcommand{\lecture}[5]{
	\pagestyle{headings}
	\thispagestyle{plain}
	\newpage
	\setcounter{chapter}{#1}
	\setcounter{page}{#2}
	\noindent
	\begin{center}
		\framebox{
			\vbox{
				\hbox to 6.28in { {\bf Algorithms
						\hfill Spring Semester, 2026} }
				\vspace{4mm}
				\hbox to 6.28in { {\Large \hfill Homework 3 (Graded) \hfill} }
				\vspace{2mm}
				\hbox to 6.28in { {\it Professor: #4 \hfill Deadline: Noon Mar 9, 2026.} }
			}
		}
	\end{center}
	\vspace*{4mm}
}

\newcommand{\topic}[2]{\section{#1} \index{#2} \markright{#1}}
\newcommand{\subtopic}[2]{\subsection{#1} \index{#2}}
\newcommand{\subsubtopic}[2]{\subsubsection{#1} \index{#2}}

\renewcommand{\cite}[1]{[#1]}
\renewcommand{\thefootnote}{\arabic{footnote}}

\newcommand{\bi}{\begin{itemize}}
	\newcommand{\ei}{\end{itemize}}
\newcommand{\be}{\begin{enumerate}}
	\newcommand{\ee}{\end{enumerate}}
\newcommand{\blank}{\vspace{1ex}}

\newtheorem{theorem}{Theorem}[chapter]
\newtheorem{lemma}[theorem]{Lemma}
\newtheorem{claim}[theorem]{Claim}
\newtheorem{digress}[theorem]{Digression}
\newtheorem{corollary}[theorem]{Corollary}

\newenvironment{dfn}{{\vspace*{1ex} \noindent \bf Definition }}{\vspace*{1ex}}
\newcommand{\bigdef}[2]{\index{#1}\begin{dfn} {\rm #2} \end{dfn}}
\newcommand{\smalldef}[1]{\index{#1} {\em #1}}

\begin{document}
	\lecture{0}{1}{\today}{Dr. Mudassir Shabbir}{180 mins.}

    \textbf{Instructions.} This is a graded homework. Submit your solutions as a single PDF. Show all work; correct answers without justification will receive no credit. You may collaborate on high-level ideas but must write up solutions independently.

\bigskip

\begin{enumerate}[label=\arabic*.]

% ---------------------------------------------------------------
% PROBLEM 1 — UNION-FIND TRACE  [10 points]
% ---------------------------------------------------------------

\item \textbf{[10 points]} \textit{Union-Find models dynamic connectivity, which arises in network routing, image segmentation, and Kruskal's MST algorithm.}

Recall the \textbf{Union-Find} data structure with two optimizations:
\begin{itemize}
    \item \textbf{Union by rank:} always attach the root of the lower-rank tree under the root of the higher-rank tree; if ranks are equal, choose either root and increment its rank by 1.
    \item \textbf{Path compression:} during \texttt{Find}$(x)$, make every node on the path from $x$ to the root point directly to the root.
\end{itemize}

Start with $n = 8$ elements $\{1, 2, 3, 4, 5, 6, 7, 8\}$, each initially its own component with rank 0.

% \begin{enumerate}[label=(\alph*)]
%     \item \textbf{[10 pts]} 
    Perform the following operations in order using union by rank.
    After \emph{each} Union, draw the resulting forest.  For every node shown, label its \textbf{rank} and indicate the \textbf{parent pointer} (or mark it as a root).
    \[
      \texttt{Union}(1,2),\quad
      \texttt{Union}(3,4),\quad
      \texttt{Union}(5,6),\quad
      \texttt{Union}(7,8)
    \]
    \[
      \texttt{Union}(1,3),\quad
      \texttt{Union}(5,7)
    \]
    \[
      \texttt{Union}(1,5)
    \]
    Then perform \texttt{Find}$(8)$ without path compression.
    Write down the sequence of nodes visited while traversing parent pointers from $8$ to the root.
    Draw the resulting forest.
    % How many distinct components exist after all operations?  List every element in each component.

%    \item \textbf{[6 pts]} Suppose that instead of union by rank you used \textbf{naive union} (always attach the second argument's root under the first argument's root).  Perform only the first four unions:
%    \[
%      \texttt{Union}(1,2),\quad \texttt{Union}(1,3),\quad \texttt{Union}(1,4),\quad \texttt{Union}(1,5)
%    \]
%    Draw the resulting tree.  What is the height of this tree, and why is this bad for subsequent \texttt{Find} operations?
% \end{enumerate}

\bigskip

% ---------------------------------------------------------------
% PROBLEM 2 — QUEUE FROM TWO STACKS  [20 points]
% ---------------------------------------------------------------

\item \textbf{[20 points]} \textit{Implementing a queue from two stacks is a classic application of amortized analysis.}

Implement a \textbf{FIFO queue} using two stacks $S_{\text{in}}$ and $S_{\text{out}}$, each supporting \texttt{Push} and \texttt{Pop} in $O(1)$ worst case, as follows:
\begin{itemize}
    \item \texttt{Enqueue}$(x)$: push $x$ onto $S_{\text{in}}$.
    \item \texttt{Dequeue}: if $S_{\text{out}}$ is non-empty, pop and return its top.
    Otherwise, move \emph{all} elements from $S_{\text{in}}$ to $S_{\text{out}}$ (by repeated pop-then-push), then pop and return the top of $S_{\text{out}}$.
\end{itemize}

\begin{enumerate}[label=(\alph*)]
    \item \textbf{[10 pts]} Trace the data structure on the following sequence of operations.
    After each operation write the contents of $S_{\text{in}}$ and $S_{\text{out}}$ (top element listed first) and the value returned (if any).

    \medskip
    \begin{center}
    \begin{tabular}{l|l|l|l}
    \hline
    \textbf{Operation} & \textbf{$S_{\text{in}}$ (top first)} & \textbf{$S_{\text{out}}$ (top first)} & \textbf{Returned} \\
    \hline
    \texttt{Enqueue}(1) & & & --- \\[4pt]
    \texttt{Enqueue}(2) & & & --- \\[4pt]
    \texttt{Enqueue}(3) & & & --- \\[4pt]
    \texttt{Dequeue}    & & & \\[4pt]
    \texttt{Enqueue}(4) & & & --- \\[4pt]
    \texttt{Dequeue}    & & & \\[4pt]
    \texttt{Dequeue}    & & & \\[4pt]
    \hline
    \end{tabular}
    \end{center}

    \item \textbf{[10 pts]} % Explain with an example why a single \texttt{Dequeue} on a queue holding $n$ elements can cost $\Theta(n)$ in the worst case.
    Use the \textbf{accounting method} to prove that any sequence of $n$ \texttt{Enqueue} and \texttt{Dequeue} operations costs $O(n)$ total.
    \footnote{Hint: Charge \textbf{3 credits} per \texttt{Enqueue} and \textbf{1 credit} per \texttt{Dequeue}, and maintain the \textbf{credit invariant}: every element currently in $S_{\text{in}}$ holds exactly 2 banked credits.}

    % For each of the three cases below, show that the credits charged are sufficient to pay the actual cost \emph{and} that the credit invariant is maintained afterwards.  Then conclude that the credit balance never goes negative, so the total actual cost is at most the total credits charged, i.e.\ $O(n)$.

    % \begin{enumerate}[label=(\roman*)]
    %     \item \texttt{Enqueue}$(x)$: actual cost $= 1$.  Where do the 3 charged credits go?

    %     \item \texttt{Dequeue} when $S_{\text{out}}$ is \textbf{non-empty}: actual cost $= 1$.  Show the 1 credit covers this; is the invariant on $S_{\text{in}}$ affected?

    %     \item \texttt{Dequeue} when $S_{\text{out}}$ is \textbf{empty} and $S_{\text{in}}$ holds $k$ elements: actual cost $= 2k + 1$.  Show that the 1 charged credit plus the 2 banked credits per element exactly cover this cost.
    % \end{enumerate}
\end{enumerate}

% ---------------------------------------------------------------
% PROBLEM 3 — KRUSKAL'S ALGORITHM: TRACE AND ROBUSTNESS  [20 points]
% ---------------------------------------------------------------

\bigskip
\item \textbf{[20 points]} \textit{Running Kruskal's algorithm by hand builds intuition for the cut property and reveals how sensitive the MST is to edge weight changes.}

Consider the following weighted undirected graph $G$ with vertices $\{1,2,3,4,5,6\}$ and edges as shown.

\begin{center}
\begin{tikzpicture}[scale=1.1,
  nd/.style={circle, draw=black, thick, minimum size=0.75cm, font=\bfseries\small},
  lbl/.style={draw=none, fill=white, inner sep=1.5pt, font=\small}]
  \node[nd] (1) at (0,   2)  {1};
  \node[nd] (2) at (2,   3)  {2};
  \node[nd] (3) at (2,   1)  {3};
  \node[nd] (4) at (4,   3)  {4};
  \node[nd] (5) at (4,   1)  {5};
  \node[nd] (6) at (6,   2)  {6};

  \draw[thick] (1) -- node[lbl, above left]  {4}  (2);
  \draw[thick] (1) -- node[lbl, below left]  {2}  (3);
  \draw[thick] (2) -- node[lbl, left]        {1}  (3);
  \draw[thick] (2) -- node[lbl, above]       {5}  (4);
  \draw[thick] (3) -- node[lbl, above]       {8}  (4);
  \draw[thick] (3) -- node[lbl, below]       {10} (5);
  \draw[thick] (4) -- node[lbl, left]        {2}  (5);
  \draw[thick] (4) -- node[lbl, above right] {6}  (6);
  \draw[thick] (5) -- node[lbl, below right] {3}  (6);
\end{tikzpicture}
\end{center}

\begin{enumerate}[label=(\alph*)]

    \item \textbf{[10 pts]} List all 9 edges in non-decreasing order of weight (break ties by smaller endpoint first), then trace Kruskal's algorithm.  Fill in every row of the table.  In the ``Action'' column write \textbf{Add} or \textbf{Skip}.  In the ``Reason'' column, for a \textbf{Skip} name the cycle that would have been created; for an \textbf{Add} state which two components merge.

    \bigskip
    \begin{center}
    \renewcommand{\arraystretch}{1.6}
    \begin{tabular}{c|c|c|c|l}
    \hline
    \textbf{Step} & \textbf{Edge} & \textbf{Weight} & \textbf{Action} & \quad\textbf{Reason / Components after} \\
    \hline
    1 & & & & \\
    2 & & & & \\
    3 & & & & \\
    4 & & & & \\
    5 & & & & \\
    6 & & & & \\
    7 & & & & \\
    8 & & & & \\
    9 & & & & \\
    \hline
    \end{tabular}
    \end{center}

    % \item \textbf{[4 pts]} Draw the MST of $G$ (label all edges and their weights) and state its total weight.

    \item \textbf{[10 pts]} When Kruskal's adds edge $(2, 4)$, it merges two components $S$ and $V \setminus S$.
    Identify the cut $(S, V \setminus S)$ at that moment, list all edges crossing it with their weights, and verify using the \textbf{cut property} that $(2,4)$ must belong to every MST of $G$.

    Then suppose the weight of $(2,4)$ is changed to $w'$.  Find the exact threshold below which the MST is unchanged.  For $w' = 9$, identify which edge replaces $(2,4)$, draw the new MST, and state its total weight.

\end{enumerate}

\end{enumerate}

\end{document}
